\documentclass[10pt,a4paper]{amsart}
\usepackage[margin=19mm]{geometry}
\usepackage{amsmath,amssymb,mathtools}
\usepackage[scr=rsfs,cal=euler]{mathalfa}
\usepackage{xfrac}
\usepackage[unicode,hidelinks,bookmarks=false]{hyperref}
\newcommand{\ie}{i.e.\ }
\newcommand{\cf}{cf.\ }
\newcommand{\resp}{respectively\ }
\newcommand{\Subsection}{\subsection}
\providecommand{\nobreakdash}{\nobreak\hbox{-}}
\setlength{\emergencystretch}{2em}
\multlinegap0pt
\usepackage{algorithm}\usepackage{algpseudocode}
\usepackage{amssymb}
\usepackage[shortlabels]{enumitem}
\usepackage[normalem]{ulem}
\usepackage{mathtools}
\usepackage{bm}
\usepackage{tikz}
\usepackage{float}
\newtheorem{lemma}{Lemma}
\title{Graph Multivector Persistence: A Unified Framework for Dynamic Systems}
\author{Donald Woukeng}
\date{Source: arXiv:2603.00677v1 (CC BY 4.0)}
\begin{document}
\maketitle

\noindent\emph{Source and licence.} Graph Multivector Persistence: A Unified Framework for Dynamic Systems. Version arXiv:2603.00677v1 (CC BY 4.0), distributed by arXiv under the Creative Commons Attribution 4.0 International licence, http://creativecommons.org/licenses/by/4.0/, as recorded in the arXiv licence metadata for this exact version and shown on the abstract page https://arxiv.org/abs/2603.00677v1. Full licence notice: \texttt{licenses/arxiv-2603.00677-CC-BY-4.0.txt}.\par\smallskip

\noindent\textit{Editorial scope.} Counted unit: the article's lemma lem:filtration\_stability (source0.tex lines 2050--2067). The article's complete original proof of it is delivered with it (source0.tex lines 2069--2118, one \texttt{proof} environment of 50 lines). No mathematics is written, repaired or re-derived: the statement and the proof are the article's own lines, in the article's own notation, and no retained line is substituted or re-flowed.  No further source-local context is needed: every cross-reference of every retained line resolves inside the two delivered spans.  Nothing was omitted from any retained span, so the package carries no omission record.  No retained line contains a citation, so the wrapper carries no bibliography and no bibliography entry.  No retained line contains a figure environment, an image-inclusion command or drawing code, so no image is delivered and the package declares no asset dependency.  Editorial formatting only, in addition: an AMS \texttt{amsart} preamble that restates the article's own declarations and macros used by the retained lines, namely the environments \texttt{lemma} on separate counters, together with the standard packages those lines need.  The retained environments print in this document as Lemma 1 in order of appearance, whereas the article prints the same items with the numbering of its own full document.  The complete native source of the article as distributed in the arXiv e-print for this exact version is bundled unchanged as \texttt{sources/195498/source0.tex}.
threshold activation of relations within $\varepsilon$ of the boundary.\begin{lemma}[Filtration stability]
\label{lem:filtration_stability}
Let $W$ and $W'$ be two relation matrices on the same finite set
with $\|W - W'\|_\infty \le \varepsilon$.  
Then, for every threshold $\lambda \in [0,1]$, the corresponding
multivector field partitions satisfy
\[
\mathcal V_{\lambda+\varepsilon}(W)
  \preceq
\mathcal V_{\lambda}(W')
  \preceq
\mathcal V_{\lambda-\varepsilon}(W),
\]
where $\mathcal A \preceq \mathcal B$ means that
$\mathcal A$ is a \emph{refinement} of $\mathcal B$
(every block of $\mathcal A$ is contained in some block of $\mathcal B$).
Equivalently, as $\lambda$ decreases, the partition becomes coarser.
\end{lemma}

\begin{proof}
Recall that the algorithm depends only on the thresholded relation
matrix
\[
w_\lambda(i,j) = \mathbf{1}_{\{W(i,j) > \lambda\}},
\]
and that lowering $\lambda$ activates additional entries (more edges
become ``active''), producing coarser partitions.

Now fix $\lambda$ and assume $\|W - W'\|_\infty \le \varepsilon$.
This means that for every pair $(i,j)$,
\[
|W(i,j) - W'(i,j)| \le \varepsilon.
\]

\textbf{Step 1.}  
Suppose $W(i,j) > \lambda + \varepsilon$.  
Then automatically $W'(i,j) \ge W(i,j) - \varepsilon > \lambda$.  
Hence every pair $(i,j)$ that is ``active'' at level $\lambda + \varepsilon$
for $W$ remains active at level $\lambda$ for $W'$.  
Therefore, every merge that occurs in $\mathcal V_{\lambda+\varepsilon}(W)$
must also occur in $\mathcal V_{\lambda}(W')$.
This implies
\[
\mathcal V_{\lambda+\varepsilon}(W) \preceq \mathcal V_{\lambda}(W').
\]

\textbf{Step 2.}  
Conversely, if $W'(i,j) > \lambda$, then
$W(i,j) \ge W'(i,j) - \varepsilon > \lambda - \varepsilon$,
so all relations active in $\mathcal V_{\lambda}(W')$ are also active in
$\mathcal V_{\lambda - \varepsilon}(W)$.
Thus the merges realized in $\mathcal V_{\lambda}(W')$ also appear in
$\mathcal V_{\lambda - \varepsilon}(W)$, giving
\[
\mathcal V_{\lambda}(W') \preceq \mathcal V_{\lambda - \varepsilon}(W).
\]

Combining the two inclusions, we obtain
\[
\mathcal V_{\lambda+\varepsilon}(W)
  \preceq
\mathcal V_{\lambda}(W')
  \preceq
\mathcal V_{\lambda-\varepsilon}(W).
\]
Since lowering $\lambda$ always coarsens the partition, this result
captures the precise way in which perturbations of $W$ translate into
a small horizontal shift of the entire filtration by at most $\varepsilon$.
\end{proof}

\end{document}
